\section*{Chapitre \ref{ch:g10:why-clouds-form} --- Pourquoi les nuages se forment}

\begin{solution}{exo:g10:why-clouds-form:1}
L'\omterm{def:g10:why-clouds-form:humidity}{humidité relative} est la teneur de l'air en vapeur d'eau divisée par sa
teneur à \omterm{def:g10:why-clouds-form:humidity}{saturation} à la même température et à la même pression, exprimée
en pourcentage. Le \omterm{def:g10:why-clouds-form:dew-point}{point de rosée} est la température à laquelle il faut
refroidir l'air, en conservant sa teneur en eau et sa pression, pour qu'il
devienne \omterm{def:g10:why-clouds-form:humidity}{saturé}.
\end{solution}

\begin{solution}{exo:g10:why-clouds-form:2}
\num{3.8}, \num{7.6}, \num{14.7} et \qty{27.2}{g/kg}. Chaque
\qty{10}{\degreeCelsius} la multiplie par environ 2 ($7.6/3.8 = 2.0$,
$14.7/7.6 \approx 1.9$, $27.2/14.7 \approx 1.9$).
\end{solution}

\begin{solution}{exo:g10:why-clouds-form:3}
Teneur à \omterm{def:g10:why-clouds-form:humidity}{saturation} à \qty{24}{\degreeCelsius}~: \qty{18.9}{g/kg}~;
\omterm{def:g10:why-clouds-form:humidity}{humidité relative} $10.0/18.9 \approx \qty{53}{\percent}$. La teneur à
\omterm{def:g10:why-clouds-form:humidity}{saturation} vaut \qty{10.0}{g/kg} à \qty{14}{\degreeCelsius}~: c'est le \omterm{def:g10:why-clouds-form:dew-point}{point
de rosée}.
\end{solution}

\begin{solution}{exo:g10:why-clouds-form:4}
En été, l'air est humide et son \omterm{def:g10:why-clouds-form:dew-point}{point de rosée} est bien supérieur à la
température du verre~: l'air qui touche le verre est refroidi sous son \omterm{def:g10:why-clouds-form:dew-point}{point
de rosée} et l'eau se condense dessus. L'air très sec d'un intérieur en hiver
a un \omterm{def:g10:why-clouds-form:dew-point}{point de rosée} inférieur à la température du verre~: refroidi à son
contact, l'air n'est toujours pas \omterm{def:g10:why-clouds-form:humidity}{saturé}.
\end{solution}

\begin{solution}{exo:g10:why-clouds-form:5}
Le chauffage par le bas (cumulus d'après-midi au-dessus d'un sol ensoleillé)~;
les montagnes (vent forcé de monter sur une chaîne)~; les fronts (air chaud
qui monte par-dessus de l'air froid)~; la convergence (vents qui confluent
près de l'équateur ou vers une dépression).
\end{solution}

\begin{solution}{exo:g10:why-clouds-form:6}
$125 \times (30 - 18) = \qty{1500}{m}$~; $125 \times (22 - 20) =
\qty{250}{m}$. La seconde journée~: les nuages se forment à seulement
\qty{250}{m} au-dessus du sol.
\end{solution}

\begin{solution}{exo:g10:why-clouds-form:7}
À \qty{1}{km}~: $18 - 9.8 = \qty{8.2}{\degreeCelsius}$~; à \qty{2}{km}~:
$18 - 19.6 = \qty{-1.6}{\degreeCelsius}$. En redescendant, elle se
réchauffe au même taux et retrouve au sol ses \qty{18}{\degreeCelsius}.
\end{solution}

\begin{solution}{exo:g10:why-clouds-form:8}
Avec \qty{4}{\degreeCelsius/km}, une parcelle non saturée
(\qty{9.8}{\degreeCelsius/km}) comme une parcelle saturée
(\qty{5}{\degreeCelsius/km}) se refroidissent plus vite que leur
environnement~: l'air est \omterm{def:g10:why-clouds-form:stability}{stable} dans les deux cas. Avec
\qty{12}{\degreeCelsius/km}, même une bulle non saturée se refroidit plus
lentement que son environnement~: elle reste plus chaude et continue de
monter, et les bulles d'air chaud s'élèvent librement.
\end{solution}

\begin{solution}{exo:g10:why-clouds-form:9}
$14.7 - 7.6 = \qty{7.1}{g}$ par kilogramme d'air. La \omterm{def:g10:earth-energy-budget:latent-heat}{chaleur latente}
réchauffe l'air dans lequel l'eau se condense, ce qui ralentit son
refroidissement.
\end{solution}

\begin{solution}{exo:g10:why-clouds-form:10}
Les volumes varient comme le cube du diamètre~: $(0.5/0.02)^3 = 25^3 = 15\,625$
gouttelettes pour une goutte de bruine~; $(2/0.02)^3 = 100^3 = 1\,000\,000$
pour une goutte de pluie.
\end{solution}

\begin{solution}{exo:g10:why-clouds-form:11}
Sous un ciel clair, le sol rayonne de l'infrarouge et se refroidit, ainsi
que l'air à son contact. Cet air froid est plus dense que l'air au-dessus~;
sur les pentes, il s'écoule vers le bas et s'accumule au fond de la vallée.
Là, il continue de se refroidir pendant la nuit jusqu'à atteindre son \omterm{def:g10:why-clouds-form:dew-point}{point
de rosée}, et des gouttelettes se forment dans l'air lui-même~: c'est le
brouillard. Sur les sommets, l'air n'est pas aussi froid et reste au-dessus
de son \omterm{def:g10:why-clouds-form:dew-point}{point de rosée}.
\end{solution}

\begin{solution}{exo:g10:why-clouds-form:12}
À l'altitude $z$, la température vaut $T - a z$ et le \omterm{def:g10:why-clouds-form:dew-point}{point de rosée}
$T_d - b z$~; ils sont égaux quand $(a - b) z = T - T_d$, donc
$z = (T - T_d)/(a - b)$. Avec $b = 1.7$~: $a - b = 8.1$, et
$1000/8.1 \approx \qty{123}{m}$ par degré.
\end{solution}

\begin{solution}{exo:g10:why-clouds-form:13}
\omterm{def:g10:why-clouds-form:condensation-level}{Niveau de condensation}~: $125 \times (15 - 11) = \qty{500}{m}$~; température
à ce niveau~: $15 - 9.8 \times 0.5 \approx \qty{10.1}{\degreeCelsius}$.
Crête~: $10.1 - 5
\times 1.5 = \qty{2.6}{\degreeCelsius}$. Descente de \qty{2}{km} au gradient
sec~: $2.6 + 9.8 \times 2 \approx \qty{22.2}{\degreeCelsius}$, environ sept
degrés de plus que du côté au vent.
\end{solution}

\begin{solution}{exo:g10:why-clouds-form:14}
À \qty{-20}{\degreeCelsius}~: $126/103 \approx 1.22$, soit
\qty{22}{\percent} de sursaturation pour la glace~; à
\qty{-10}{\degreeCelsius}~: $287/260 \approx
1.10$, soit \qty{10}{\percent}. Entre ces températures, plus le nuage est
froid, plus l'avantage de la glace est grand~: le processus fonctionne le
mieux dans la partie haute et froide d'un nuage, là où gouttelettes et
cristaux sont mêlés.
\end{solution}

\begin{solution}{exo:g10:why-clouds-form:15}
Non saturée~: la parcelle se refroidit de \qty{9.8}{\degreeCelsius}, son
environnement de \qty{6.5}{\degreeCelsius}~: elle est plus froide de
\qty{3.3}{\degreeCelsius} et redescend. Saturée~: elle se refroidit de
\qty{5}{\degreeCelsius}, elle est donc plus chaude de
\qty{1.5}{\degreeCelsius} et continue de monter. L'air n'est \omterm{def:g10:why-clouds-form:stability}{instable} qu'à
une condition~: que l'air qui monte soit \omterm{def:g10:why-clouds-form:humidity}{saturé}.
\end{solution}

\begin{solution}{pb:g10:why-clouds-form:1}
\textbf{1.} \qty{14.7}{g/kg}.

\textbf{2.} $11.4/14.7 \approx \qty{78}{\percent}$.

\textbf{3.} \qty{16}{\degreeCelsius}, où la teneur à \omterm{def:g10:why-clouds-form:humidity}{saturation} vaut
\qty{11.4}{g/kg}.

\textbf{4.} $11.4/24.1 \approx \qty{47}{\percent}$.

\textbf{5.} Non~: il contient les mêmes \qty{11.4}{g/kg}. Son \omterm{def:g10:why-clouds-form:humidity}{humidité
relative} a baissé parce que, plus chaud, il pouvait en contenir davantage.

\textbf{6.} À son \omterm{def:g10:why-clouds-form:dew-point}{point de rosée}, \qty{16}{\degreeCelsius}.

\textbf{7.} $11.4 - 8.7 = \qty{2.7}{g}$ par kilogramme d'air.

\textbf{8.} Sous un ciel nocturne clair, l'herbe et les toits rayonnent de
l'infrarouge et deviennent plus froids que l'air~; l'air qui les touche est
refroidi sous son \omterm{def:g10:why-clouds-form:dew-point}{point de rosée} avant le reste de l'air.

\textbf{9.} $(28.0 - 13.4)/1.5 \approx \qty{9.7}{\degreeCelsius/km}$,
pratiquement le gradient sec~: les bulles de l'après-midi ont brassé cette
couche.

\textbf{10.} $(13.4 - 0.4)/2.0 = \qty{6.5}{\degreeCelsius/km}$.

\textbf{11.} \omterm{def:g10:why-clouds-form:stability}{Stable} pour une parcelle non saturée, qui se refroidit plus
vite (\qty{9.8}{\degreeCelsius/km})~; \omterm{def:g10:why-clouds-form:stability}{instable} pour une parcelle saturée,
qui se refroidit plus lentement (\qty{5}{\degreeCelsius/km}).

\textbf{12.} La température augmente avec l'altitude, de \num{-0.3} à
\qty{1.2}{\degreeCelsius}~: une couche très \omterm{def:g10:why-clouds-form:stability}{stable}, qui joue le rôle de
couvercle. Le nuage qui monte s'y arrête~: ses sommets sont vers
\qty{4}{km}.

\textbf{13.} $28 - 16 = \qty{12}{\degreeCelsius}$.

\textbf{14.} Température $28 - 9.8 = \qty{18.2}{\degreeCelsius}$~; \omterm{def:g10:why-clouds-form:dew-point}{point de
rosée} $16 - 1.8 = \qty{14.2}{\degreeCelsius}$.

\textbf{15.} De 12 à $18.2 - 14.2 = 4.0$~: il s'est réduit de
\qty{8}{\degreeCelsius}, soit un rythme de \qty{8}{\degreeCelsius} par
kilomètre.

\textbf{16.} Un degré d'écart est comblé en $1/8$ km, soit \qty{125}{m}.

\textbf{17.} $125 \times (20 - 16) = \qty{500}{m}$.

\textbf{18.} Le \omterm{def:g10:why-clouds-form:dew-point}{point de rosée} reste à \qty{16}{\degreeCelsius} tandis que
l'air proche du sol se réchauffe~: l'écart grandit, et avec lui l'altitude
à laquelle il se comble.

\textbf{19.} $28 - 9.8\,z = 16 - 1.8\,z$.

\textbf{20.} $8\,z = 12$, donc $z = 1.5$~: la base des cumulus de
l'après-midi est \textbf{à environ \qty{1.5}{km}} au-dessus de la plaine (le
calcul exact avec les propriétés de l'eau donne \qty{1.51}{km}).
\end{solution}
